\section*{Conclusion}

Reporting delays are commonly understood as shifting and smoothing epidemic signals. Our results show that their more consequential effect is to impose a quantitative, frequency-dependent limit on what can be inferred from downstream surveillance data. By combining the transfer function of the delay process with conversion and observation-model-specific uncertainty, we developed a spectral separability framework that characterizes whether competing upstream epidemic trajectories can be statistically distinguished after they pass through the observation process. This perspective reframes epidemic reconstruction from recovering a latent trajectory to determining which differences between epidemic histories, and which scientific conclusions derived from them, remain supported by the information transmitted downstream.

Across empirical delay distributions, this framework revealed a consistent hierarchy across temporal scales and inferential targets. Slow epidemic variation is comparatively well preserved, whereas short-term fluctuations e.g. weekly and sub-weekly fluctuations can be attenuated by orders of magnitude and may require unrealistically large upstream contrasts to remain distinguishable from noise. This loss of information does not affect all derived quantities equally. Features governed primarily by broad, low-frequency structure, such as the timing of a major epidemic wave, can remain comparatively stable across admissible reconstructions, while quantities based on local temporal differences, such as short-term growth rates, depend more strongly on weakly identified high-frequency components. The same principle extends to event attribution. Even when the occurrence and timing of a lockdown, mass gathering, or vaccination campaign are known, its transmission magnitude and temporal footprint may occupy an extended low-separability region rather than correspond to a uniquely identifiable parameterization. Identifiability is therefore not a binary property of an entire epidemic trajectory. It is specific to the temporal scale, epidemiological feature, and scientific comparison being considered.

These limits clarify the role of regularization, prior distributions, and mechanistic structure in epidemic reconstruction. Such assumptions are indispensable for stabilizing an ill-posed inverse problem and for producing epidemiologically coherent estimates. They cannot, however, recreate information that has been strongly suppressed by the observation process. When the likelihood provides little discrimination among alternative high-frequency trajectories, different priors, smoothing penalties, process models, or preprocessing choices may select different solutions while achieving nearly indistinguishable downstream fits. A reconstruction may consequently be numerically stable, have narrow conditional uncertainty, and closely reproduce the observed data without its fine-scale features being strongly identified by those data. Downstream goodness of fit, inferential precision, and statistical identification should therefore be treated as distinct properties. 

The quantitative boundary between identifiable and weakly identifiable variation is conditional on the observation process through which epidemic information is transmitted. Among its components, uncertainty in the delay distribution is particularly consequential because the delay kernel directly determines the frequency-dependent attenuation of upstream variation. Our sensitivity analysis shows that moderate changes in the
assumed median delay or dispersion can shift the identifiability boundary by orders of magnitude for rapid perturbations, whereas longer-timescale variation is generally less sensitive to the same uncertainty (Appendix~\ref{sec:Appendix_S6}; Fig.~\ref{fig:delay-sensitivity}). Consequently, the temporal resolution limit should be evaluated across epidemiologically plausible delay specifications rather than under a single fixed kernel,
particularly when inference concerns temporal scales near the resolution boundary.

Other components of the observation process—including ascertainment, observation noise, and the observation window—likewise influence the quantitative location of this boundary. Complementary observations, such as wastewater measurements, testing data, mobility or contact indicators, genomic information, and administrative records of interventions, can further reduce ambiguity by contributing information unavailable from a
single delayed outcome stream. Their contribution should be interpreted as providing additional sources of identification, rather than as evidence that the original surveillance series itself contained the missing temporal information.

These findings have implications for both surveillance design and the evaluation of epidemic models. Shortening and narrowing reporting delays, increasing ascertainment, reducing measurement noise, and combining observations with complementary temporal responses do more than improve timeliness: they increase the range of epidemic variation that can be statistically resolved. Model evaluation should ask not only whether an estimate fits observed outcomes or produces a precise posterior, but also whether the particular epidemiological quantity being reported is identifiable under the assumed observation process and remains stable across plausible delay specifications and reconstruction assumptions. Making these limits explicit would shift epidemic inference away from presenting a single detailed latent history toward reporting which temporal scales, epidemiological features, and scientific conclusions are actually supported by the available surveillance system.